% ============================================================
% Flash-Sale Microservices Capstone — Report Skeleton
% Student: Vu Dinh Kiet (523V0011) — Class 23K50201 — TDTU
% Week 1 skeleton per FLASHSALE_EXECUTION_CHECKLIST.md Step 1.4.
%
% DRAFT NOTICE: Introduction, Problem Statement, Objectives, Scope
% and Limitations below are first-draft text written from the
% checklist's own quoted excerpts of the proposal (de-cuong-cap-nhat.md /
% the LaTeX proposal), NOT from the proposal's original wording, because
% that source file is not present in this repo/session. Per Step 1.4's
% own warning, do not leave this as a silent copy of the proposal's
% pitch-register text either way -- reread against your actual proposal
% and reword into report register (documents what WAS built/found,
% not what SHOULD be approved) before this is considered final.
% ============================================================

\documentclass[12pt,a4paper]{report}

\usepackage[utf8]{inputenc}
\usepackage[T5]{fontenc} % Vietnamese text support; swap for polyglossia+xelatex if compiling with XeLaTeX/LuaLaTeX instead
\usepackage{geometry}
\geometry{margin=2.5cm}
\usepackage{graphicx}
\usepackage{hyperref}
\usepackage{listings}
\usepackage{booktabs}
\usepackage{amsmath}
\usepackage{caption}

% NOTE on diagrams: Mermaid blocks used throughout docs/*.md (state machine,
% ERD, sequence diagrams) do NOT render directly in LaTeX like they do in
% Markdown/GitHub. Export each one as an image (mermaid.live -> PNG/SVG, or
% the mermaid-cli `mmdc`) into report/figures/ and \includegraphics it where
% noted below. Budget time for this conversion step before each week's
% report-writing task, not just at the end.

\title{Thiết kế và Đánh giá hiệu năng hệ thống Microservices \\
xử lý giao dịch song song quy mô lớn trong kịch bản Flash-sale}
\author{Vũ Đình Kiệt (523V0011) \\ Class 23K50201 \\ Ton Duc Thang University}
\date{\today}

\begin{document}

\maketitle
\tableofcontents

% ============================================================
\chapter{Introduction}
% ============================================================

Flash-sale events -- short, fixed-duration promotional windows in which a large
volume of buyers compete for a much smaller pool of inventory -- are a deliberately
adversarial workload for backend systems: request volume spikes by orders of
magnitude in seconds, and correctness (never selling stock that does not exist)
must hold even as throughput far exceeds normal operating conditions.

This project designs, implements, and experimentally evaluates a microservices-based
order-and-inventory system built specifically to survive this workload shape. Rather
than treating "handle high traffic" as a single monolithic goal, the work isolates
two mechanisms that are conflated in most informal Flash-sale discussions: how the
system stays \emph{correct} under concurrency (Section~\ref{sec:problem}), and how it
stays \emph{responsive} under overload it cannot fully absorb (load shedding,
Week 12). Both are evaluated through controlled, repeatable JMeter-driven experiments
against multiple architectural configurations (C0 through the final resilient
design), not through a single end-to-end demo.

% ============================================================
\chapter{Problem Statement}
\label{sec:problem}
% ============================================================

Under concurrent access to a shared, finite inventory count, a naively implemented
order path exhibits a classic read-then-write race: two concurrent requests can each
observe available stock, both decide independently that a reservation is valid, and
both commit -- overselling the product. This is demonstrated directly and
deliberately in this project (the C0 baseline, Week 3) before being solved.

Beyond single-request correctness, a Flash-sale system built as a set of
independently-deployed services introduces a second class of problem: partial
failure. A service can crash between updating its own state and notifying another
service of that update (the dual-write problem, addressed by the Transactional
Outbox pattern, Week 6); a message can be delivered more than once by the broker
(addressed by idempotent consumers and the Inbox pattern, Week 9); and the system as
a whole can receive more concurrent requests than it is provisioned to process
correctly, in which case naive behavior degrades unpredictably rather than
predictably (addressed by load shedding, Week 12).

The problem this project addresses, precisely: \emph{design and build an
order-and-inventory subsystem that provably avoids overselling under concurrent
load, remains correct despite partial service failure and at-least-once message
delivery, and degrades predictably rather than catastrophically once load exceeds
provisioned capacity -- then measure, with controlled experiments, what that
correctness and resilience cost in throughput and latency compared to a simpler
baseline.}

% ============================================================
\chapter{Objectives}
% ============================================================

\begin{enumerate}
  \item Implement and empirically compare at least two inventory-reservation
    strategies -- a synchronous PostgreSQL atomic-update baseline (C1) and an
    asynchronous, event-driven Redis-backed design -- under identical experimental
    conditions.
  \item Guarantee, and experimentally validate, the four correctness invariants
    defined in \texttt{docs/00-scope-lock.md}: bounded reservations, controlled
    overload response, replica scalability, and traceable state.
  \item Implement the resilience mechanisms a production event-driven system
    requires -- Transactional Outbox, idempotent/Inbox-pattern consumers, bounded
    retry with a Dead-Letter Queue, and a reconciliation safety net -- and
    demonstrate each one specifically defeats the failure mode it targets.
  \item Measure system behavior (throughput, p50/p95/p99 latency, correctness
    violation rate) across the full range from normal load through deliberate
    overload, and report the point at which -- and the mechanism by which -- the
    system sheds load rather than failing uncontrolled.
\end{enumerate}

% ============================================================
\chapter{Scope}
% ============================================================

See \texttt{docs/00-scope-lock.md} for the authoritative, frozen scope statement.
Summary for report readers:

\textbf{In scope:} Order Service and Inventory Service (ABP Framework,
app-nolayers), a Process Worker simulating downstream processing, PostgreSQL
(schema-per-service) and Redis (Lua-scripted atomic reservation) for state,
RabbitMQ for asynchronous messaging with Transactional Outbox and Inbox patterns,
Nginx-based load shedding, and a JMeter-driven experimental protocol comparing
multiple configurations (C0 naive, C1 synchronous baseline, and the full
asynchronous design) under normal and overload traffic.

\textbf{Out of scope:} frontend UI, product catalog/cart/promotions, user
management and authentication/authorization, and any clustered/replicated
deployment of Redis, RabbitMQ, or PostgreSQL (single-instance only). Payment
processing and failure-compensation logic beyond the Process Worker's deterministic
simulated success path are explicitly excluded.

% ============================================================
\chapter{Limitations}
% ============================================================

The following are acknowledged, deliberate limitations of the experimental design
and system, not oversights to correct later within this project's timeline:

\begin{itemize}
  \item Single-instance infrastructure (Redis, RabbitMQ, PostgreSQL) means this
    project does not evaluate behavior under infrastructure-level failure
    (node loss, network partition, replica failover) -- only application- and
    service-level failure modes are covered.
  \item The Process Worker simulates downstream processing with a deterministic
    delay and a successful outcome; it does not model realistic downstream failure
    rates, so the \texttt{ProcessingFailed} state's real-world frequency is not
    experimentally characterized (see the open question in
    \texttt{docs/order-state-machine.md}).
  \item Experiments run in a single local/cloud environment (Week 13's AWS EC2
    deployment) rather than across a distributed/multi-region topology, so
    network-latency effects representative of real multi-region Flash-sale traffic
    are out of scope.
  \item Load-shedding thresholds (Nginx rate limits, application-level concurrency
    caps) are tuned empirically against this project's specific hardware and
    workload shape; they are not claimed to generalize to arbitrary deployment
    scale without re-tuning.
  \item \texttt{GET /api/orders/\{id\}} (Week 5) performs no ownership check --
    any caller who has an order's \texttt{id} can read it. This is the direct
    consequence of excluding authentication/authorization from scope (this
    document's own Scope chapter), not an oversight: order IDs are
    cryptographically random GUIDs, functioning as an unguessable capability
    token rather than a sequential identifier an unrelated party could
    enumerate. A production system would still pair this with real
    authorization; this project does not, on purpose. Client-facing
    idempotency keys (\texttt{POST /api/orders}) are hardened against the
    adjacent risk of key collision instead -- reusing a key with a different
    request body is rejected (\texttt{409}) rather than silently returning an
    unrelated order, so at minimum two unrelated callers colliding on the same
    key surface as an explicit conflict, not a silent data leak.
\end{itemize}

% ============================================================
\chapter{System Requirements}
% ============================================================

\section{Functional Requirements}
\begin{itemize}
  \item Accept an order request for a single product/quantity and return
    immediately with a pending state, without blocking on inventory confirmation.
  \item Allow a client to query an order's current state and full transition
    history at any time after submission.
  \item Reserve inventory such that concurrent requests for the same product can
    never together reserve more than the available stock.
  \item Detect and reject duplicate order submissions (client-level) and duplicate
    message deliveries (broker-level) without double-processing either.
  \item Recover automatically from a crashed consumer, a temporarily unreachable
    broker, or a message that repeatedly fails processing, without operator
    intervention for the common cases.
  \item Shed excess load predictably (explicit rejection) once request rate
    exceeds provisioned capacity, rather than degrading unpredictably.
\end{itemize}

\section{Non-Functional Requirements}
\begin{itemize}
  \item \textbf{Correctness under concurrency} takes priority over latency — see
    the four guarantees in \texttt{docs/00-scope-lock.md}.
  \item \textbf{Observability}: every order's lifecycle must be reconstructable
    end-to-end from logs via a single correlation ID (Week 11).
  \item \textbf{Horizontal scalability}: no service may depend on in-memory-only
    state that would break when running more than one replica.
  \item \textbf{Reproducibility}: every experiment run must be preceded by a
    deterministic reset/seed step (\texttt{scripts/reset-and-seed.ps1}, Week 3) so
    results are comparable across configurations and across repeated runs.
\end{itemize}

% ============================================================
\chapter{Technology Selection}
% ============================================================

\begin{description}
  \item[PostgreSQL 16] Chosen for the durable, transactional side of the system
    (orders, outbox, inbox) specifically because the Transactional Outbox pattern
    (Week 6) requires the business write and the event write to commit atomically
    in one transaction — a relational engine with real ACID transactions is a
    prerequisite for that pattern, not just a convenient default.
  \item[Redis Stack] Chosen over plain Redis specifically for its bundled Lua
    scripting support: the inventory reservation check-and-decrement
    (\texttt{scripts/redis/reserve.lua}, Week 7) must execute as one atomic
    operation server-side, which Lua's single-threaded execution inside Redis
    guarantees without needing external locking.
  \item[RabbitMQ vs. Kafka] RabbitMQ was chosen over Kafka for this workload
    shape: the system needs per-message routing, retry, and dead-lettering
    semantics (Week 10) that map directly onto RabbitMQ's exchange/queue/DLX
    model, whereas this project doesn't need Kafka's log-retention/replay
    strengths (no consumer ever needs to replay historical events from scratch).
  \item[ABP Framework (\texttt{app-nolayers})] Chosen for built-in event bus
    abstractions (\texttt{Volo.Abp.EventBus.RabbitMQ}) and EF Core integration
    that reduce boilerplate around the parts of the system that aren't the actual
    research question (reservation correctness, resilience, load shedding). The
    \texttt{app-nolayers} template specifically — not the default layered
    template — to keep each service's codebase proportional to its actual
    complexity for a project this scoped.
  \item[Nginx] Chosen as the load-shedding entry point (Week 12) for its mature,
    well-documented \texttt{limit\_req} rate-limiting module, which is sufficient
    for this project's single-node local/EC2 deployment without needing a full
    API gateway.
\end{description}

% ============================================================
\chapter{Development Environment}
% ============================================================

\section{Local infrastructure (Docker Compose)}

\texttt{infra/docker-compose.yml} defines four local dependencies, brought up with
\texttt{docker compose up -d} from \texttt{infra/}:

\begin{center}
\begin{tabular}{@{}lll@{}}
\toprule
Service & Image & Ports \\
\midrule
PostgreSQL & \texttt{postgres:16} & 5432 \\
Redis Stack & \texttt{redis/redis-stack:latest} & 6379 (data), 8001 (RedisInsight UI) \\
RabbitMQ & \texttt{rabbitmq:3-management} & 5672 (AMQP), 15672 (management UI) \\
Nginx & \texttt{nginx:latest} & 8080 (host) $\to$ 80 (container) \\
\bottomrule
\end{tabular}
\end{center}

Dev credentials (\texttt{flashsale} / \texttt{flashsale\_dev}) are local-only —
never reused for the Week 13 AWS EC2 deployment, which uses separate
Terraform-managed credentials.

\section{Toolchain}

Verified working versions as of Week 2 setup (see \texttt{setup.ps1} for the
onboarding script that installs/checks all of these on a fresh machine):

\begin{center}
\begin{tabular}{@{}ll@{}}
\toprule
Tool & Version \\
\midrule
Git & 2.53 \\
.NET SDK & 10.0.300 \\
ABP CLI & 10.5.0 \\
Docker Desktop & 29.4.3 \\
Terraform & 1.15.8 \\
Apache JMeter & 5.6.3 \\
\bottomrule
\end{tabular}
\end{center}

Builds run from native Windows PowerShell against the \texttt{D:\textbackslash}
drive, not inside a WSL2 shell — cross-filesystem I/O between Windows and WSL2 is
the main source of slow \texttt{docker build}/bind-mount performance on this setup
(see \texttt{FLASHSALE\_EXECUTION\_CHECKLIST.md}, Step 0.2).

% ============================================================
\chapter{System Analysis and Design}
% ============================================================

\section{Entity-Relationship Diagram}
% TODO Week 2 -- \includegraphics figure exported from docs/erd.md (schema-per-service: order_service, inventory_service).

\section{API Contract}

ABP generates the OpenAPI/Swagger contract directly from the running
controllers, rather than it being hand-written -- the full document is
exported to \texttt{docs/api-contract-v1.json} (Week 4) and re-exported as
each week's endpoints change, so it can never drift from the actual code the
way a hand-maintained spec could.

Reviewing the Week 4 export against \texttt{docs/order-state-machine.md}
surfaced a scope-fidelity issue worth recording rather than quietly fixing:
of roughly 30 exposed endpoint groups, only \texttt{/api/orders} (the C1
baseline from Week 3) actually belongs to this system. Every other path
(\texttt{/api/identity/*}, \texttt{/api/account/*},
\texttt{/api/permission-management/*}, \texttt{/api/multi-tenancy/*}, ...)
is default scaffolding from ABP modules the project's own scope lock
(\texttt{docs/00-scope-lock.md}) explicitly excludes -- user management,
authentication/authorization, and multi-tenancy. None of it traces to any
state-machine transition, which is exactly what Step 4.1's review criterion
is designed to catch. These modules were pulled in automatically by the
\texttt{app-nolayers} template's default package set (Week 3.1) and were not
a deliberate choice; trimming them is a reasonable Week 5 cleanup item, since
touching module wiring isn't itself Week 4 scope.

The real API surface -- idempotent order creation and status lookup -- is
built in Week 5; this section is re-exported and expanded once that exists.

\section{Event Contract}

Full contract: \texttt{docs/event-contract.md}; class definitions:
\texttt{src/FlashSale.EventContracts/}. Three events, one per fully-specified
state-machine transition (\texttt{docs/order-state-machine.md}):
\texttt{OrderPlacedEto} ($[*] \to$ \texttt{PendingStock}),
\texttt{StockReservedEto} (\texttt{PendingStock} $\to$ \texttt{Confirmed}),
and \texttt{StockRejectedEto} (\texttt{PendingStock} $\to$ \texttt{Rejected}).
Each carries a \texttt{MessageId} -- the broker-redelivery idempotency key
used by the Week 9 Inbox pattern, distinct from the client-facing
\texttt{Idempotency-Key} header built in Week 5.

The contract deliberately omits a \texttt{CorrelationId} field and a
\texttt{Reason} field on the rejection event. Both are easy to imagine
wanting eventually, but nothing produces or consumes either one yet:
correlation-ID propagation is explicitly a Week 11 addition, and the only
rejection cause currently in scope (insufficient stock) gives a
\texttt{Reason} field nothing to distinguish between. Both are documented as
deliberate omissions rather than oversights, precisely so a later reviewer
(or a later week's version of the author) doesn't mistake "not needed yet"
for "forgotten."

The contract lives in a separate class library
(\texttt{FlashSale.EventContracts}) rather than inside Order Service itself,
referenced by whichever service needs it. This is what makes "the contract
doc and the code can't drift" actually true across service boundaries, not
just within one service: Inventory Service (Week 7) will reference the same
compiled types Order Service publishes, rather than maintaining its own copy
that could silently diverge.

\section{Order State Machine}
See \texttt{docs/order-state-machine.md} for the full diagram and
transition table. Two transitions
(\texttt{Confirmed} $\to$ \texttt{Processing} $\to$
\texttt{Completed}/\texttt{ProcessingFailed}) remain an explicitly recorded
open question pending the Week 11 Process Worker design -- inventing a
trigger for them now, to make the diagram look more finished than the design
actually is, would cost more credibility at defense time than leaving the
gap visible.

\section{Sequence Diagrams}
Full diagrams: \texttt{docs/sequence-diagrams.md} (happy path and rejection
path). Originally drafted in Week 4 as target design; revised in Week 8
(Step 8.4) against what was actually built and verified live through that
week -- the Outbox on both services, both publisher workers, and the
Redis-backed reservation round trip. Both diagrams still stop short of the
\texttt{Processing} state for the same reason the state machine does above
-- extending a sequence diagram past a transition whose trigger isn't known
yet means drawing a mechanism that doesn't exist.

\section{Design Rationale}
The central design decision this project makes -- comparing a synchronous
PostgreSQL-atomic-update baseline (C1, Week 3) against an asynchronous,
event-driven design (Weeks 5--12) under the \emph{same} correctness
invariants (\texttt{docs/inventory-invariants.md}) -- is what the
Experimental Design and Results chapters ultimately quantify. Every
mechanism introduced from here forward (Outbox, Inbox, DLQ, load shedding)
exists to answer one question precisely: what does buying resilience and
horizontal scalability actually cost in throughput and latency, relative to
the simpler baseline that cannot provide either.

% ============================================================
\chapter{Implementation}
% ============================================================

\section{Baseline Configurations — C0 and C1 (Week 3)}
\label{sec:c0-c1}

Before building the real asynchronous design (Weeks 5--12), two controlled
baseline configurations were built and compared, entirely inside the Order
Service, with no separate Inventory Service involved yet. Both share a single
raw-SQL table, \texttt{order\_service.inventory} (deliberately \emph{not}
EF Core-migrated -- see \texttt{docs/erd.md} -- since C0 and C1 are baseline
scaffolding, not the final design), seeded to a known stock value before every
run via \texttt{scripts/reset-and-seed.ps1}.

\subsection{C0 -- the naive baseline (a controlled failure demonstration)}

C0 (\texttt{src/FlashSale.OrderService/Experiments/C0Naive/C0NaiveDemo.cs})
implements the textbook race condition: read current stock, check it in
application code, then issue a separate, unconditional \texttt{UPDATE} --
two round trips, no locking between them:

\begin{verbatim}
SELECT stock FROM order_service.inventory WHERE product_id = @p;
-- (application code checks: stock > 0 ?)
UPDATE order_service.inventory SET stock = stock - 1 WHERE product_id = @p;
\end{verbatim}

C0 is deliberately not reachable via any HTTP route -- it is invoked only via
\texttt{dotnet run -{}- -{}-run-c0-demo}, which bypasses the whole ABP host.
This is a stronger isolation guarantee than "don't add a controller for it":
there is no route in the routing table at all for a stray request to hit
later in the project.

Five concurrent attempts were fired against a product seeded with
\texttt{stock = 1}, synchronized on a shared start gate so every attempt
reaches the read at the same instant (without this, the race did not
reproduce on every run -- scheduling jitter alone was occasionally enough to
serialize the two "concurrent" reads). The result:

\begin{verbatim}
[C0] Seeded 'c0-demo-product' with stock = 1. Firing 5 concurrent naive requests...
[C0] Attempt 3: saw stock=1, Confirmed.
[C0] Attempt 1: saw stock=1, Confirmed.
[C0] Attempt 4: saw stock=1, Confirmed.
[C0] Attempt 2: saw stock=1, Confirmed.
[C0] Attempt 5: saw stock=1, Confirmed.
[C0] Result: 5 of 5 concurrent requests got 'Confirmed' for a product seeded with stock = 1.
[C0] OVER-SOLD: more than one request reserved the same single unit of stock.
[C0] Final stock in database: -4 (started at 1).
\end{verbatim}

All five requests believed they had validly reserved the single unit of
stock; the database was left at \(-4\). This is the exact failure mode C1
exists to eliminate. Note that \texttt{order\_service.inventory} deliberately
has no \texttt{CHECK (stock >= 0)} constraint -- an earlier version of this
table had one, and it only turned the bug into an unhandled exception
instead of letting it manifest, which is a form of debugging the bug away
rather than capturing it.

\subsection{C1 -- the PostgreSQL atomic baseline}

C1 (\texttt{src/FlashSale.OrderService/Controllers/BaselineOrdersController.cs})
replaces the two round trips with a single conditional statement, where the
\texttt{WHERE} clause is evaluated by PostgreSQL against the row's current
value under its own row lock, not against a value the application read
earlier and might now be stale:

\begin{verbatim}
UPDATE order_service.inventory
SET stock = stock - 1
WHERE product_id = @productId AND stock >= 1
RETURNING stock;
\end{verbatim}

This runs inside the same database transaction as the order-outcome log
insert, exposed as a real synchronous endpoint, \texttt{POST /api/c1/orders}
(moved off \texttt{/api/orders} in Week 5 once that path became the real
design's endpoint -- Week 13/14 need both configurations independently
reachable for comparison, not one path whose behavior depends on a hidden
flag). The same five-concurrent-request test used against C0 was repeated against
C1, seeded with \texttt{stock = 1}: \textbf{exactly 1 request Confirmed, 4
Rejected, final stock 0} -- never negative, never more than one winner.
\texttt{SELECT ... FOR UPDATE} followed by a separate \texttt{UPDATE} was
deliberately avoided here: that reintroduces C0's exact race window between
the two statements unless wrapped very carefully, whereas a single
\texttt{UPDATE...WHERE...RETURNING} is simpler and provably atomic.

\subsection{Correctness criteria}

The four invariants defined in \texttt{docs/inventory-invariants.md} apply
directly to this comparison: C0 violates Invariant~1
(\texttt{available\_inventory >= 0}) and Invariant~2
(\texttt{successful\_reservations <= initial\_inventory}) by construction
under concurrency; C1 satisfies both, verified empirically above rather than
only argued from the SQL's atomicity.

\subsection{JMeter smoke test}

A 10-thread, single-iteration JMeter plan
(\texttt{tests/jmeter/smoke-test.jmx}) targets C1's \texttt{POST /api/c1/orders}
against a product seeded with ample stock (\texttt{flash-product-1},
\texttt{stock = 1000}), confirming the full HTTP pipeline works end-to-end
before any load testing begins. The first run of this test returned a clean
\texttt{404} on all 10 requests -- worth recording precisely because it
\emph{looked} like a pass against the checklist's own "2xx/4xx, not a
connection error" criterion, but wasn't: the saved test plan had no port set,
so JMeter defaulted to port 80, which on this machine belongs to an unrelated
process (WSL2's \texttt{wslrelay.exe}), not the Order Service (which also
wasn't running at the time). Fixed by setting the port explicitly and
re-running against the actual service:

\begin{center}
\begin{tabular}{@{}ll@{}}
\toprule
Metric & Result \\
\midrule
Requests & 10 \\
Responses & 10 $\times$ \texttt{200 OK} \\
Errors & 0 \\
Avg / Min / Max latency & 913\,ms / 480\,ms / 1240\,ms \\
\texttt{flash-product-1} stock & 1000 $\to$ 990 (exactly 10 decrements) \\
\bottomrule
\end{tabular}
\end{center}

Raw results: \texttt{tests/jmeter/results/smoke-test-result.jtl}.

\section{Order Service}

\subsection{API design}

Two endpoints (\texttt{Controllers/OrdersController.cs}):
\texttt{POST /api/orders} (idempotent creation) and
\texttt{GET /api/orders/\{id\}} (status lookup). Both return immediately from
Postgres state -- there is no event publication yet (that arrives with the
Transactional Outbox, Week 6), so every order created this week stays in
\texttt{PendingStock} forever, which is expected for this week's scope, not
a defect.

\subsection{Idempotency strategy}

A required \texttt{Idempotency-Key} request header is checked against
existing orders before creating a new one; a repeated key returns the
original order unchanged rather than creating a duplicate. This is the
\emph{client}-facing idempotency layer -- entirely distinct from the
\emph{broker}-message idempotency the Week 9 Inbox pattern adds later, which
dedupes RabbitMQ redelivering the same message rather than a client
resubmitting the same request. Conflating the two was flagged as a risk as
early as \texttt{docs/00-scope-lock.md} (Week 1); keeping the two mechanisms'
names and purposes distinct in the code (\texttt{IdempotencyKey} the client
supplies vs.\ \texttt{MessageId} on the ETOs in
\texttt{FlashSale.EventContracts}) is what keeps that distinction from
blurring once both exist side by side.

Verified manually against the running service (not just by inspection):
sending the same \texttt{Idempotency-Key} twice returns the identical order
\texttt{id} both times, and a direct Postgres query confirms exactly one row
exists for that key -- the second request did not insert a duplicate.

\subsection{State management}

\texttt{Entities/Order.cs} implements the state machine from
\texttt{docs/order-state-machine.md} as guarded transition logic
(\texttt{Order.TransitionTo}), not an unguarded status field: an illegal
transition throws \texttt{InvalidOrderStateTransitionException} rather than
silently succeeding. Per that document's own open question, \texttt{Confirmed}
currently has no legal outgoing transition either -- \texttt{Confirmed ->
Processing}'s trigger depends on the Week 11 Process Worker design, so no
code path can legally perform it yet. The alternative (guessing a trigger to
make the state machine look more finished) would cost more at defense time
than an honestly-incomplete transition table.

\subsection{Test cases}

14 unit tests (\texttt{tests/FlashSale.OrderService.Tests}), all passing:
every legal transition succeeds and stamps the expected timestamp; every
state currently without a legal outgoing transition (including
\texttt{Confirmed}, precisely because of the open question above) throws on
every attempted target; and the exception itself is asserted to report the
correct from/to states. C0 and C1 (Week 3) are deliberately excluded from
this test suite -- they are one-shot demonstrations captured as report
evidence, not code with a future correctness guarantee to protect.

\begin{verbatim}
Passed!  - Failed: 0, Passed: 14, Skipped: 0, Total: 14
\end{verbatim}

\subsection{Event Publication Reliability (Week 6) -- the dual-write problem and the Transactional Outbox}
\label{sec:outbox}

\paragraph{The problem.} Once order creation needs to notify another service
(Inventory Service, Week 7) that stock should be reserved, the order write
and the "notify" action become two separate operations against two separate
systems (Postgres and RabbitMQ) that cannot be wrapped in one atomic
transaction the way two SQL statements against the same database can. Naively
sequencing them --- write the order, then publish an event --- exposes two
distinct failure orderings, and a system that only handles one of them is
still broken:

\begin{enumerate}
  \item \textbf{The database commits but the publish fails.} The order now
    exists in \texttt{PendingStock} with no corresponding event ever sent.
    Nothing will ever move it out of \texttt{PendingStock} -- it is silently
    stuck forever, and neither the client nor any operator has any signal
    that anything is wrong.
  \item \textbf{The publish succeeds but the database transaction rolls
    back} (or the process crashes between the two steps). Inventory Service
    now receives an \texttt{OrderPlaced} event for an order that, as far as
    Order Service's own database is concerned, never happened. It may
    reserve real stock against a phantom order that no client will ever be
    able to query or complete.
\end{enumerate}

Both failures share a root cause: two systems that must agree with each
other are being updated as if they were one, with no mechanism enforcing
that they actually end up agreeing.

\paragraph{How the Outbox pattern eliminates both.} Rather than publishing
directly, the order-creation handler writes two rows to \emph{one}
database -- the \texttt{orders} row and a row in \texttt{order\_service.outbox\_events}
recording that an \texttt{OrderPlaced} event needs to be published -- inside
the exact same \texttt{SaveChangesAsync} call
(\texttt{Controllers/OrdersController.cs}). Because both rows are ordinary
rows in the same relational database, Postgres's own transaction guarantee
applies to both simultaneously: either both commit, or neither does. This
directly closes both failure orderings above --- there is no longer a window
where the order exists without a corresponding pending event, or vice versa,
because "the order exists" and "an event needs publishing" have been reduced
to a single fact instead of two facts that could disagree.

This claim is verified directly, not just argued from the pattern's
description, in
\texttt{tests/FlashSale.OrderService.Tests/OutboxAtomicityTests.cs}: a real
SQLite in-memory database (not a mock, and not EF Core's InMemory provider,
which does not enforce relational constraints and so cannot demonstrate a
rollback at all) is made to reject an outbox insert via a genuine primary-key
collision, and the accompanying order insert is confirmed absent afterward
--- the whole transaction rolled back, not just the failing half.

\begin{verbatim}
Failed_outbox_insert_rolls_back_the_order_insert_too: PASS
Order_and_outbox_row_both_commit_in_one_SaveChanges_call: PASS
\end{verbatim}

\paragraph{Getting the pending event to RabbitMQ.} Writing the outbox row
only records that publication is needed -- a separate background worker
(\texttt{Messaging/OutboxPublisherWorker.cs}) polls for unpublished rows
every two seconds and publishes each one, marking a row published only once
the broker has actually confirmed receipt, not merely once the publish call
returned without throwing. This distinction matters: a broker that accepts a
TCP write and then rejects the message asynchronously would let the weaker
guarantee silently drop it. RabbitMQ.Client 7.x's
\texttt{CreateChannelOptions(publisherConfirmationsEnabled: true,
publisherConfirmationTrackingEnabled: true)} makes \texttt{BasicPublishAsync}
itself not return until the broker has acknowledged the message, which is
what \texttt{OutboxEvent.MarkPublished} is called on the strength of.

\texttt{Volo.Abp.EventBus.RabbitMQ}'s own options
(\texttt{AbpRabbitMqOptions}/\texttt{AbpRabbitMqEventBusOptions}) were
checked, per this step's own instruction not to assume a configuration
property exists without verifying it, and do not expose a documented way to
enable publisher confirms. Publishing is done via the raw
\texttt{RabbitMQ.Client} package instead, for exactly this reason -- reading
connection settings from the same JSON shape ABP's own package uses, for
consistency, without depending on a whole framework module for a
configuration convenience.

\paragraph{Failure handling in the worker itself.} A publish attempt that
fails is retried up to three additional times with exponential backoff
(1s/2s/4s) before the worker gives up for that poll cycle and leaves the row
unpublished for the next one -- this was verified against a real outage, not
simulated: stopping the RabbitMQ container mid-test produced exactly this
log sequence, with the service continuing to accept new orders throughout
(the worker's own tick failures are caught and logged, never allowed to
crash the host):

\begin{verbatim}
[WRN] Failed to publish outbox event ... (attempt 1/4); retrying in 00:00:01
[WRN] Failed to publish outbox event ... (attempt 2/4); retrying in 00:00:02
[WRN] Failed to publish outbox event ... (attempt 3/4); retrying in 00:00:04
[ERR] Failed to publish outbox event ... after 4 attempts; leaving unpublished,
      will retry on the next poll
\end{verbatim}

Restarting the broker let the same row publish successfully on a later poll
without any manual intervention, confirming the retry loop is not just
present in the code but actually recovers.

\paragraph{Verifying delivery, not just the publish call succeeding.} Since
Inventory Service does not exist until Week 7, there is no consumer yet to
observe the message being consumed -- the checklist anticipates exactly this
("stay at 1 if you haven't built the consumer yet, which is expected until
Week 7"). A queue (\texttt{OrderPlaced}) was bound to the exchange
(\texttt{flashsale.order.exchange}) as a temporary stand-in verification
target via RabbitMQ's HTTP management API, and each order created during
testing produced exactly one message that persisted in that queue (message
count increased by exactly one per order, including the one that survived
the broker outage above) -- confirming messages are actually reaching
RabbitMQ, not just that the publish call returns successfully. This queue is
verification scaffolding, not a permanent design element; Week 7's real
Inventory Service consumer declares and binds its own durable queue in code.

\section{Inventory Service}

\subsection{Scaffolding, done differently from Order Service on purpose}

Inventory Service was scaffolded with a deliberately minimal ABP module set
(\texttt{MVC}, \texttt{Autofac}, \texttt{EntityFrameworkCore.PostgreSql},
\texttt{Swashbuckle}, \texttt{Serilog} only) rather than the full default set
\texttt{abp new} produces. Order Service (Week 3) carried the default set
(Identity, Account, PermissionManagement, TenantManagement,
FeatureManagement, SettingManagement, OpenIddict, AuditLogging, the LeptonX
UI theme), and that turned out to cost more than unused weight: it directly
caused the Week 3 migration/schema bug (Setting Management's runtime queries
don't resolve a non-default schema consistently) and, per Week 4's API
contract review, produced roughly 29 of 30 exposed endpoints that traced to
nothing in the state machine. Since \texttt{docs/00-scope-lock.md} already
excludes user management, auth/authz, and multi-tenancy entirely, Inventory
Service simply never adds those modules, rather than adding and later
removing them. The concrete result: its Swagger contract exposes exactly the
three unavoidable ABP core meta-endpoints and nothing else, versus Order
Service's roughly thirty endpoint groups. This is a deliberate difference
between the two services' scaffolds, not an inconsistency -- Order Service's
equivalent cleanup remains a flagged, not-yet-done item, tracked separately
rather than retrofitted here.

\subsection{Redis data structure and the Lua reservation script}

Actual stock lives in Redis, not Postgres (\texttt{docs/erd.md}): an integer
key \texttt{inventory:\{productId\}} for available stock, and a set key
\texttt{processed:\{productId\}} recording which order IDs have already been
reserved for that product -- the second key is what makes the duplicate
check possible without touching Postgres at all.

\begin{verbatim}
if redis.call('SISMEMBER', KEYS[2], ARGV[1]) == 1 then
  return 'DUPLICATE'
end
local stock = tonumber(redis.call('GET', KEYS[1]))
if stock == nil or stock <= 0 then
  return 'REJECTED'
end
redis.call('DECR', KEYS[1])
redis.call('SADD', KEYS[2], ARGV[1])
return 'RESERVED'
\end{verbatim}

Redis executes a Lua script as one atomic, single-threaded operation --
this is what makes the duplicate-check, stock-check, decrement, and
mark-processed sequence race-free without any application-level locking.
The three outcomes were verified directly against \texttt{redis-cli}, before
any C\# integration existed, isolating the script's own correctness from
the .NET client's:

\begin{verbatim}
> SET inventory:test 1
> EVAL <script> 2 inventory:test processed:test order-1
RESERVED
> EVAL <script> 2 inventory:test processed:test order-1
DUPLICATE
> EVAL <script> 2 inventory:test processed:test order-2
REJECTED
> GET inventory:test
0
\end{verbatim}

All three outcomes matched expectations exactly, and stock never went
negative or was double-decremented. The C\# side
(\texttt{Inventory/InventoryReservationService.cs}) is a thin wrapper that
loads this exact script file at startup and passes \texttt{KEYS}/\texttt{ARGV}
through -- it does not, and must not, reimplement any of the check-then-act
logic in C\#, since doing so would reintroduce the exact race condition the
Lua script's atomicity exists to prevent.

\subsection{Inventory warm-up}

\texttt{scripts/warm-up.ps1} loads known starting stock for each test
product before a Flash-sale "opens". Product availability is conditional on
this having run by construction, not by convention: an unset inventory key
reads as Lua's \texttt{nil}, which the script already treats as
\texttt{REJECTED}, so a product nobody warmed up simply cannot be reserved.

\subsection{OrderPlaced handling}

\texttt{Messaging/OrderPlacedConsumer.cs} consumes via the raw
\texttt{RabbitMQ.Client}, the same choice and for the same reason as Order
Service's \texttt{OutboxPublisherWorker} (Week 6): messages are published as
plain JSON with a routing key, not through \texttt{Volo.Abp.EventBus.RabbitMQ}'s
own envelope format, so a consumer built against that abstraction would not
be looking at the same message shape. Keeping publish and consume on the
same primitives end to end avoids that mismatch. Each message is
deserialized into \texttt{OrderPlacedEto} (\texttt{FlashSale.EventContracts},
Week 4), reserved via the Lua script above, and manually acknowledged only
after successful processing -- a failure nacks with requeue, and a
transient connection failure retries the whole consumer connection with a
5-second backoff rather than crashing the service. Week 9 replaces this
simple ack-after-success strategy with the full Inbox-pattern-based
idempotent acknowledgement.

Publishing \texttt{StockReserved}/\texttt{StockRejected} back to Order
Service is Week 8's addition (Section~\ref{sec:week8}); this week's own
verification below stops at the Lua/consumer boundary, which is sufficient
to confirm the reservation logic itself end to end before wiring the result
back out.

\paragraph{End-to-end verification.} Rather than only checking Redis state
after the fact, the full path was watched live: an order created via
\texttt{POST /api/orders} was observed flowing through the Outbox, into
RabbitMQ, into this consumer, into the Lua script, and into an updated Redis
key, within about two seconds end to end (RabbitMQ pushes to consumers; this
is not a polling delay). Both outcomes were exercised through the complete
pipeline, not just at the Lua/redis-cli level above:

\begin{verbatim}
[INF] OrderPlaced 7918a432-... for product flash-product-1 -> Reserved
[INF] OrderPlaced 1dd50a8a-... for product out-of-stock-product -> Rejected
\end{verbatim}

A useful side effect of building the consumer this week: three
\texttt{OrderPlaced} messages left over from Week 6's manual verification
queue (including one that had survived a deliberate RabbitMQ outage) were
still sitting unconsumed. Starting this consumer drained all three
automatically on first connection, decrementing \texttt{flash-product-1}'s
stock from 1000 to 997 and recording exactly three entries in
\texttt{processed:flash-product-1} -- confirming not only that new messages
flow through correctly, but that messages published before any consumer
existed are not lost, exactly what the Outbox pattern from Week 6 was
supposed to guarantee.

\section{Closing the Loop -- Stock Results and Order State (Week 8)}
\label{sec:week8}

Weeks 6--7 built one direction of the pipeline: an order becomes an
\texttt{OrderPlaced} event, which Inventory Service reserves against Redis.
Week 8 builds the return direction -- Inventory Service telling Order
Service what happened, and Order Service acting on it -- which is what
actually lets an order leave \texttt{PendingStock}.

\subsection{Inventory Service's own Outbox -- a harder dual-write problem}

Order Service's Week 6 Outbox writes a business row and an outbox row to the
\emph{same} Postgres database in one transaction, so atomicity is free.
Inventory Service cannot do that: the business fact (stock was reserved) is
recorded in Redis by \texttt{reserve.lua}, and the outbox fact (tell Order
Service) must be recorded in Postgres -- two different data stores that
cannot share a transaction. Rather than accept a gap where a process crash
between the two writes silently drops the notification, Inventory Service
gets its own \texttt{outbox\_events} table (\texttt{inventory\_service}
schema) with one deliberate difference from Order Service's: the unique
index is on \texttt{OrderId}, not \texttt{Id}.

That difference is what makes the Lua script's \texttt{DUPLICATE} outcome
safe. \texttt{DUPLICATE} means a \emph{previous} delivery of the same
\texttt{OrderPlaced} message already reserved the stock in Redis -- so if
the consumer treated it as a no-op, a redelivered message would never
produce a \texttt{StockReserved} row, and that order would sit in
\texttt{PendingStock} forever despite the reservation having genuinely
succeeded. Instead, both \texttt{RESERVED} and \texttt{DUPLICATE} run the
same \texttt{EnsureResultOutboxEventAsync} step: check whether an outbox row
for this \texttt{OrderId} already exists, and only insert if it doesn't. The
unique index turns a possible race between two redeliveries into a caught
\texttt{DbUpdateException} rather than a duplicate row -- "ensure exists",
not "insert unconditionally". \texttt{RabbitMqOutboxPublisher} and
\texttt{OutboxPublisherWorker} on the Inventory Service side are otherwise
the same design as Order Service's Week 6 versions (raw
\texttt{RabbitMQ.Client}, publisher confirms, 2s poll interval).

\subsection{Order Service: consuming the result}

\texttt{Messaging/StockResultConsumer.cs} binds two queues,
\texttt{StockReserved} and \texttt{StockRejected}, to
\texttt{flashsale.order.exchange}. Both share one handler,
\texttt{HandleResultAsync}, parameterized by the target state: look up the
order; if it is already in the target state, acknowledge as a no-op
(handles redelivery without re-running the transition); otherwise call
\texttt{Order.TransitionTo}. A transition rejected by the state machine
(\texttt{InvalidOrderStateTransitionException} -- e.g. a message claims
\texttt{Reserved} for an order already \texttt{Rejected}) is a genuine
conflict, not a redelivery, and is logged loudly and nacked \emph{without}
requeue, since there is no DLQ yet (Week 10) to catch a message that would
otherwise requeue forever. Any other exception nacks \emph{with} requeue.
Unlike \texttt{EnsureResultOutboxEventAsync} above,
\texttt{SaveChangesAsync()} here is called directly rather than through
ABP's \texttt{SaveChangesOnDbContextAsync} bypass -- this code runs inside a
real DI scope resolved via \texttt{scope.ServiceProvider}, where the
services ABP's own \texttt{SaveChangesAsync} depends on (audit logging,
event dispatch) are actually available; the bypass is only needed when a
\texttt{DbContext} is constructed directly outside DI, as in an isolated
unit test.

\paragraph{End-to-end verification.} Both outcomes were exercised through
the complete, real pipeline -- not a unit test double for Redis, RabbitMQ,
or Postgres. An order posted against a product with available stock was
polled from \texttt{PendingStock} through to \texttt{Confirmed}; a second
order against a product forced to zero stock was polled through to
\texttt{Rejected}. In both cases, Postgres showed exactly one Inventory
Service outbox row for the order, with the expected event type and
\texttt{Published = t}.

\subsection{Automated correctness validation (Step 8.3)}

\texttt{scripts/validate-correctness.ps1} turns the four invariants from
\texttt{docs/inventory-invariants.md} into a runnable check against the live
system, parameterized by product id and expected initial stock rather than
hardcoded -- Weeks 13--14 reuse this exact script, unmodified, against
several different product/stock configurations, so it cannot bake in one
week's test values. Two PowerShell-specific bugs surfaced and were fixed
while building it: passing SQL containing double-quoted, case-sensitive
Postgres identifiers as a \texttt{docker exec ... -c} argument silently
stripped the quotes (fixed by piping the SQL via stdin instead, matching
\texttt{reset-and-seed.ps1}'s existing pattern), and a helper function
originally named its parameter \texttt{\$args}, which collides with
PowerShell's own reserved automatic variable of the same name (renamed to
\texttt{\$redisArgs}).

Run against a clean, controlled batch (Postgres tables truncated, Redis
re-seeded, five orders created across two products, specifically so the
script's first real trial wasn't validating against leftover mixed test
data):

\begin{verbatim}
Invariant                                        Status    Detail
1. available_inventory >= 0                      PASS
2. successful_reservations <= initial_inventory  PASS
3. <=1 successful reservation per order          PASS      (structurally
                                                             guaranteed by PK;
                                                             verified, not
                                                             assumed)
4. duplicate delivery never double-deducts stock PENDING   order_service.
                                                             processed_messages
                                                             does not exist
                                                             yet (Week 9)

RESULT: All checkable invariants hold.
\end{verbatim}

Invariant 4 is reported \texttt{PENDING}, not \texttt{PASS}, because the
table it needs (\texttt{order\_service.processed\_messages}) is Week 9's
Inbox pattern and does not exist yet -- the script checks for the table's
existence first and only runs the real duplicate-message check if present.
Reporting a fake pass here would be claiming a check that never actually
ran; the script's exit code (0) reflects only the invariants that were
actually checkable this week, and \texttt{PENDING} does not fail it.

\subsection{Sequence diagrams brought current (Step 8.4)}

\texttt{docs/sequence-diagrams.md}, drafted in Week 4 as target design, was
revised this week against what was actually built and verified above: the
response code for order creation is \texttt{201 Created} (the diagram had
drawn \texttt{202 Accepted}), Inventory Service's result publication is
shown as the same two-hop Outbox-plus-publisher-worker pattern as Order
Service's own rather than a single inline publish, and the
\texttt{DUPLICATE} handling above is called out explicitly rather than left
implicit. Nginx remains drawn as target design, since routing through it
isn't exercised until Week 12.

\section{Reliability and Idempotency Layer (Week 9)}
\label{sec:week9}

RabbitMQ's own guarantee is at-least-once delivery -- a message can be, and
in the Step 9.3 test below actually was, delivered more than once. Weeks
6--8 already made every individual state write safe against that
(Order.TransitionTo has no legal self-loop, and reserve.lua's own
\texttt{DUPLICATE} result and the Week 8 outbox's unique-OrderId index both
tolerate replays). Week 9 adds the piece those didn't cover: a broker-message-
level Inbox, so a redelivered message is recognized and skipped \emph{before}
any of that per-message logic even runs, and so the report can state
precisely what turns RabbitMQ's at-least-once delivery into effectively-once
\emph{side effects} -- because those are not the same guarantee, and the
distinction is exactly what this section is required to make explicit.

\subsection{Two idempotency layers, not one (Step 9.1)}

It would be easy to assume Redis's \texttt{processed:\{productId\}} set
(Week 7) already covers this. It doesn't -- it answers a different question.
That set dedupes by \emph{order id}, at the business level: two different
messages must not both reserve stock for the same order. The new
\texttt{processed\_messages} table (one per schema, \texttt{order\_service}
and \texttt{inventory\_service}, docs/erd.md) dedupes by \emph{MessageId}, at
the broker-delivery level: the exact same physical RabbitMQ message must not
be processed twice. A bug that somehow re-published \texttt{OrderPlaced} for
an order with a fresh MessageId would be caught by the Redis-level check but
sail straight through a naive MessageId-only check, and vice versa for a
literal broker redelivery -- keeping both is not redundancy, it's covering
two different failure shapes with one constraint each. Both tables share the
same shape (\texttt{id uuid PK}, \texttt{message\_id uuid UK}, \texttt{message\_type},
\texttt{processed\_at}), applied via \texttt{dotnet ef migrations add
AddProcessedMessages} against each service, and confirmed directly:

\begin{verbatim}
order_service.processed_messages:     "MessageId" UNIQUE, btree
inventory_service.processed_messages: "MessageId" UNIQUE, btree
\end{verbatim}

\subsection{Idempotent consumer logic (Step 9.2)}

Both consumers' actual correctness logic was pulled out of the
RabbitMQ-facing classes into two small static classes --
\texttt{StockResultProcessor} (Order Service) and \texttt{OrderPlacedProcessor}
(Inventory Service) -- leaving \texttt{StockResultConsumer} and
\texttt{OrderPlacedConsumer} as pure broker plumbing (ack/nack,
deserialization) around them. Both follow the same shape: check the Inbox by
MessageId first; if present, ack as a no-op without touching the business
state at all; if absent, apply the business change and insert the Inbox row,
both tracked on the same DbContext and committed by one
\texttt{SaveChangesAsync} call, so a crash between "state changed" and
"MessageId recorded" is impossible by construction -- either both happened
in that one transaction, or neither did.

One deliberate departure from a naive "if already in target state, skip"
check: \texttt{Order.TransitionTo} has no legal self-loop for
\texttt{Confirmed}/\texttt{Rejected} (both map to an empty list in
\texttt{Order.LegalTransitions}), so calling it again on an order already in
the target state would incorrectly throw
\texttt{InvalidOrderStateTransitionException} on a perfectly legitimate
redelivery. \texttt{StockResultProcessor} branches around this explicitly --
skipping the transition call itself while still recording the Inbox row --
rather than letting that exception do double duty for both "genuine
conflict" and "safe redelivery", which would make the two indistinguishable
to the caller.

On the Inventory Service side, the Postgres-level Inbox check runs
\emph{before} the Lua script is even called -- an exact redelivery is
recognized and skipped without a Redis round trip at all, which is a real
efficiency gain, not just symmetry with the Order Service side for its own
sake.

\subsection{Acknowledgement strategy, proven with a real crash (Step 9.3)}

Both consumers already used manual acknowledgement (\texttt{autoAck: false}), acking
only after the local transaction above commits. Rather than assert this from
reading the code, it was proven directly: a temporary artificial delay was
inserted between the "message received" log line and the start of
processing (removed again before this commit -- \texttt{git diff} on this
commit shows no trace of it), an order was placed, and the moment Inventory
Service logged receipt, its process was killed outright
(\texttt{Stop-Process -Force} on the specific \texttt{FlashSale.InventoryService.exe}
PID, not the \texttt{dotnet run} wrapper around it):

\begin{verbatim}
[INF] Received OrderPlaced 70d0502d-... for order 6da814db-...
SUCCESS: The process with PID 3484 has been terminated.
\end{verbatim}

Confirmed directly, before restarting anything: Redis stock unchanged, zero
outbox rows and zero Inbox rows for this order/message, order still
\texttt{PendingStock}. The message was never acked, so it was never
removed from the queue -- exactly what manual acknowledgement is for.
Restarting Inventory Service triggered RabbitMQ to redeliver that same
unacked message (confirmed by the identical MessageId reappearing in the new
process's logs, not a new publish), and this second delivery ran to
completion:

\begin{verbatim}
[INF] Received OrderPlaced 70d0502d-... for order 6da814db-...    (redelivery)
[INF] OrderPlaced 6da814db-... for product flash-product-1 -> Reserved
\end{verbatim}

Final state, checked directly rather than inferred: Redis stock decremented
by exactly one (999 $\to$ 998, not 997), exactly one outbox row
(\texttt{Published = t}), exactly one Inbox row for that MessageId, and the
order transitioned to \texttt{Confirmed} exactly once. RabbitMQ delivered
the message twice; the system produced the business side effect once -- the
precise distinction this section opened with, now demonstrated rather than
asserted.

\subsection{Duplicate-message test and re-validation (Step 9.4)}

\texttt{DuplicateMessageIdempotencyTests.cs} drives
\texttt{StockResultProcessor} directly with the same MessageId twice (the
same SQLite-real-engine reasoning as \texttt{OutboxAtomicityTests}: a unique
constraint is the thing under test, and EF Core's InMemory provider doesn't
enforce those). Three cases: a redelivered message is a no-op and leaves
exactly one Inbox row; two \emph{different} messages for an order already in
its target state are both safely absorbed without throwing; and a genuine
conflict (a second message claiming a different terminal state for an
already-Rejected order) still throws
\texttt{InvalidOrderStateTransitionException} and leaves no Inbox row behind
for the failed attempt. \texttt{StockResultProcessor} exposes a
save-delegate-injectable overload so these tests can substitute
\texttt{SaveChangesOnDbContextAsync} for the real, DI-dependent
\texttt{SaveChangesAsync} -- the same documented limitation
\texttt{OutboxAtomicityTests} already works around, not a new one.
Full suite: 19/19 passing (16 from Weeks 5--8, 3 new).

Re-running \texttt{scripts/validate-correctness.ps1} (Step 8.3) against the
live state left over from the Step 9.3 crash test is the actual point of
Step 9.4's second instruction, and it is what finally makes Invariant 4 a
real result instead of a placeholder:

\begin{verbatim}
Invariant                                        Status Detail
1. available_inventory >= 0                      PASS   Current stock: 998
2. successful_reservations <= initial_inventory  PASS   5 <= 1000 seeded
3. <=1 successful reservation per order          PASS   (PK-guaranteed)
4. duplicate delivery never double-deducts stock PASS   No duplicate message
                                                          ids in processed_messages

RESULT: All checkable invariants hold.
\end{verbatim}

Invariant 4 was reported \texttt{PENDING} in Week 8 because
\texttt{processed\_messages} didn't exist yet; the script needed no code
change at all to go from \texttt{PENDING} to a real \texttt{PASS} -- it was
already written (Step 8.3) to check for the table's existence and run the
real query once it appeared, precisely so this week wouldn't require
touching it.

\section{Failure Handling and Recovery}
% TODO Week 10 -- bounded retry, DLQ, timeout handling, reconciliation job.

\section{Process Worker and Observability}
% TODO Week 11 -- Process Worker, correlation ID, structured logging, metrics.

\section{Load Shedding}
% TODO Week 12 -- Nginx rate limiting, application-level concurrency limiting.

% ============================================================
\chapter{Experimental Design}
% ============================================================
% TODO Week 4/13 -- docs/test-plan.md summary; configurations compared (C0/C1/final design); load profiles.

% ============================================================
\chapter{Results}
% ============================================================
% TODO Week 14 -- p50/p95/p99 tables and throughput-over-time charts per configuration (build empty templates in Week 11 per Step 11.4).

% ============================================================
\chapter{Discussion}
% ============================================================
% TODO Week 14.

% ============================================================
\chapter{Conclusion}
% ============================================================
% TODO Week 15.

% ============================================================
% \backmatter is a book-class-only command (undefined in the "report" class
% used here, docs/00-scope-lock.md's \documentclass line) -- \chapter* below
% already gives References an unnumbered chapter, and \addcontentsline
% already puts it in the TOC, so \backmatter had no effect it wasn't
% already achieving and was simply a compile-breaking leftover.
\chapter*{References}
\addcontentsline{toc}{chapter}{References}
% TODO -- bibliography (consider biblatex instead of manual list once sources accumulate).

\appendix
\chapter{Appendices}
% TODO -- full API contract JSON, full event contract, additional raw experiment data.

\end{document}
